%----------------------------------------------------------------------------------------
%	ABSTRACT PAGES
%----------------------------------------------------------------------------------------

% IMPORTANT NOTE: the abstract must always be written in two languages. If the report
% is written in Portuguese you have selected 'portuguese' as the language in the document class.
% Therefore, the portuguese version of the abstract must come first, so write it in the
% below area denoted by 'MAIN LANGUAGE ABSTRACT'. The english version follows in the
% 'SECOND LANGUAGE ABSTRACT' section.
% If the report is written in English, first will come the abstract in English
% ('MAIN LANGUAGE ABSTRACT') and then in Portuguese ('SECOND LANGUAGE ABSTRACT').

\begin{abstract}
%%%%%%%%%%%%%%%%%%%%%%%%%%%%%% MAIN LANGUAGE ABSTRACT %%%%%%%%%%%%%%%%%%%%%%%%%%%%%%%%%%
Traditional acoustic stethoscopes frequently undergo silent material degradation, such as polyvinyl chloride (PVC) hardening and eartip fatigue, which severely attenuates clinically relevant cardiopulmonary sounds. Despite this diagnostic risk, these instruments remain largely exempt from routine metrological verification. This project aimed to develop a low-cost, integrated electromechanical calibration system capable of objectively characterizing stethoscope frequency response and signal fidelity. By utilizing a surface exciter coupled to a tissue-simulating elastomer phantom, the physical prototype successfully bypasses free-air acoustic decoupling to generate controlled, structure-borne vibrations across the 20 to 2000 Hz diagnostic bandwidth. 

Data acquisition was driven by an ESP32 microcontroller sampling at 20,000 Hz, safely exceeding the Nyquist threshold, which fed a Python-based Short-Time Fourier Transform (STFT) spectral analysis pipeline. Controlled sinusoidal testing confirmed the successful capture of critical low-frequency cardiovascular targets at 20 Hz and 50 Hz, alongside high-frequency pulmonary signals up to 1111 Hz without attenuation or digital blindness. However, the spectral analysis also accurately identified system limitations, specifically non-linear amplitude compression introduced by the pre-amplifier Automatic Gain Control (AGC) and Total Harmonic Distortion (THD) at 200 Hz caused by single-supply amplifier clipping. 

Ultimately, the developed prototype successfully demonstrates a functional proof of concept for structure-borne stethoscope metrology. While analog hardware refinements are required to achieve perfect linearity, the system establishes a quantifiable, sustainable diagnostic tool designed to detect mechanical degradation and prevent compromised acoustic instruments from remaining in active clinical service.



%----------------------------------------------------------------------------------------

\vspace*{10mm} 
\noindent
\textbf{\keywordslabel}: Stethoscope Calibration, Acoustic Transducer, Digital Signal Processing 

%%%%%%%%%%%%%%%%%%%%%%%%% END OF THE MAIN LANGUAGE ABSTRACT %%%%%%%%%%%%%%%%%%%%%%%%%%%%%%
\end{abstract}
\begin{secondlangabstract}
%%%%%%%%%%%%%%%%%%%%%%%%%%%%%% SECOND LANGUAGE ABSTRACT %%%%%%%%%%%%%%%%%%%%%%%%%%%%%%%%%%
Os estetoscópios acústicos tradicionais sofrem frequentemente uma degradação material silenciosa, como o endurecimento do policloreto de vinilo (PVC) e a fadiga das olivas, o que atenua severamente os sons cardiopulmonares clinicamente relevantes. Apesar deste risco de diagnóstico, estes instrumentos permanecem, na sua maioria, isentos de verificação metrológica de rotina. Este projeto teve como objetivo desenvolver um sistema de calibração eletromecânico integrado e de baixo custo, capaz de caracterizar objetivamente a resposta em frequência e a fidelidade do sinal dos estetoscópios. Ao utilizar um excitador de superfície acoplado a um fantoma de elastómero que simula o tecido biológico, o protótipo físico contorna com sucesso o desacoplamento acústico em ar livre para gerar vibrações mecânicas estruturais controladas ao longo da largura de banda de diagnóstico de 20 a 2000 Hz. 

A aquisição de dados foi controlada por um microcontrolador ESP32 a realizar amostragens a 20.000 Hz, superando com segurança o critério de Nyquist, o qual alimentou uma linha de processamento de análise espetral baseada na Transformada de Fourier de Curto Termo (STFT) em Python. Os testes sinusoidais controlados confirmaram a captação bem-sucedida de alvos cardiovasculares críticos de baixa frequência a 20 Hz e 50 Hz, juntamente com sinais pulmonares de alta frequência até 1111 Hz sem atenuação ou cegueira digital. No entanto, a análise espetral também identificou com precisão as limitações do sistema, especificamente a compressão de amplitude não linear introduzida pelo Controlo Automático de Ganho (AGC) do pré-amplificador e a Distorção Harmónica Total (THD) a 200 Hz causada pelo ceifamento (clipping) no amplificador de alimentação simples. 

Em suma, o protótipo desenvolvido demonstra com sucesso uma prova de conceito funcional para a metrologia estrutural de estetoscópios. Embora sejam necessários refinamentos no hardware analógico para alcançar uma linearidade perfeita, o sistema estabelece uma ferramenta de diagnóstico quantificável e sustentável, concebida para detetar a degradação mecânica e evitar que instrumentos acústicos comprometidos permaneçam no ativo em contexto clínico.


%----------------------------------------------------------------------------------------

\vspace*{10mm} 
\noindent
\textbf{\keywordslabel}: Calibração de estetoscópios, transdutor acústico, processamento digital de sinais

%%%%%%%%%%%%%%%%%%%%%%%%%% END OF THE SECOND LANGUAGE ABSTRACT %%%%%%%%%%%%%%%%%%%%%%%%%%%%%
\end{secondlangabstract}
